\chapter{Extension points and cross-cutting contracts}\label{ch:extensions}
\section{Plugins and dependency injection}
Application startup creates shared services through Guice modules. The plugin mechanism discovers types marked with \code{@Plugin}, identifies plugin interfaces, and binds matching implementations into injectable sets. Consumers request the relevant service set rather than hard-coding every implementation in the window that uses it. This supports families such as extension providers and optimization services.

Discovery and execution are different stages. Discovering an extension provider allows the UI to offer a capability. An individual simulation still owns its selected extension and configuration. A plugin JAR is application code with access to the process; it is not isolated merely because it was discovered dynamically.

\sourceline{\source{core/src/main/java/info/openrocket/core/plugin/Plugin.java}{Plugin}; \source{core/src/main/java/info/openrocket/core/plugin/PluginModule.java}{PluginModule}; \source{core/src/main/java/info/openrocket/core/plugin/AnnotationFinderImpl.java}{AnnotationFinderImpl}.}

\section{Simulation extensions versus listeners}
A \class{SimulationExtension} is the persistent, configurable feature attached to a simulation. It supplies identity and metadata, stores configuration, can declare result types, participates in document loading, and initializes runtime behavior. A \class{SimulationListener} is the callback object that actually observes or modifies an execution. One extension can create several collaborating runtime objects.

\class{AbstractSimulationExtension} copies its \class{Config} on cloning and on configuration access. This protects the configuration map from accidental sharing. It does not automatically deep-copy arbitrary subclass fields, open files, static variables, or external library state. An extension intended for repeated or parallel execution must explicitly define what is shared, what is immutable, and what is freshly created.

The FC extension follows this pattern: saved settings describe the computer, while initialization creates a listener that owns the per-flight runtime, sensor adapter, control bridges, and logs. The run's copied rocket prevents actuator outputs from rewriting the editor's design. Extension configuration, the FC library artifact, and transient firmware state are three different lifetimes.

\sourceline{\source{core/src/main/java/info/openrocket/core/simulation/extension/SimulationExtension.java}{SimulationExtension}; \source{core/src/main/java/info/openrocket/core/simulation/extension/AbstractSimulationExtension.java}{AbstractSimulationExtension}; \source{core/src/main/java/info/openrocket/core/simulation/extension/impl/ZephyrusFlightComputer.java}{ZephyrusFlightComputer}.}

\section{Callback families and numerical consequences}
\begin{tabularx}{\linewidth}{@{}L{4cm}YY@{}}\toprule
Interface & Typical role & Important boundary\\\midrule
\class{SimulationListener} & Run start/end and before/after-step behavior & A completed step differs from an internal derivative evaluation\\
\class{SimulationEventListener} & Event notifications and permitted event decisions & A listener may influence whether a queued event is accepted\\
\class{SimulationComputationListener} & Observe or replace atmospheric, mass, aerodynamic, thrust, or acceleration calculations & Calculations may occur several times within one RK step\\
\class{AbstractSimulationListener} & Default implementations across the callback families & Defaults and return sentinels are part of the API\\\bottomrule
\end{tabularx}

The listener helper invokes callbacks and interprets their return values. Depending on the callback, a boolean controls continuation or acceptance, a non-null object replaces a computed value, or a non-NaN scalar supplies an override. A new listener should follow the exact method contract rather than assuming all callbacks have the same convention.

Ordering matters when two listeners modify the same quantity. A random perturbation listener, an actuator listener, and a logging listener can otherwise observe different stages of the computation. Time-varying inputs should be functions of physical time and stable run identity where possible. Drawing a fresh random value every time the RK solver asks for acceleration makes the answer depend on the solver's internal evaluation count.

The FC's acquisition contract is more restrictive than a generic acceleration callback: it captures the appropriate first derivative state, associates it with the start of an accepted interval, and publishes coherent measurements after acceptance. Its virtual scheduler and deadline limiter then coordinate task completion with the engine. A generic \code{postStep} callback by itself does not reproduce that timing model.

\sourceline{\source{core/src/main/java/info/openrocket/core/simulation/listeners/SimulationListener.java}{SimulationListener}; \source{core/src/main/java/info/openrocket/core/simulation/listeners/SimulationEventListener.java}{SimulationEventListener}; \source{core/src/main/java/info/openrocket/core/simulation/listeners/SimulationComputationListener.java}{SimulationComputationListener}; \source{core/src/main/java/info/openrocket/core/simulation/listeners/SimulationListenerHelper.java}{SimulationListenerHelper}.}

\section{Scripts, Java boards, and legacy controllers}
The inherited scripting extension evaluates a configured script through an available script engine and exposes listener callbacks. On document load it disables enabled scripts that are not locally trusted. The language must be supported by the runtime and provide the required invocable interface. A script is therefore both a persistence concern and a runtime dependency.

MIT Java-board programs use the FC's own API, compilation, source-hash approval, measurement context, and task scheduler. They are not interchangeable with arbitrary OpenRocket simulation listeners. A board sees the FC context and submits commands through controlled bridges; a general simulation listener operates at the engine callback boundary. Both run in-process, but their intended APIs and scheduling contracts differ.

The source also retains older controller/listener paths, including \class{NewControlStepListener} and \class{AirbrakesControllerListener}. Their names and presence do not mean every FC run uses them as its controller implementation. Follow the actual registration and execution path from \class{Simulation} and the selected extensions. The current default RK6 selection and legacy timestep settings are especially easy to overlook when reading upstream classes alone.

\sourceline{\source{core/src/main/java/info/openrocket/core/simulation/extension/impl/ScriptingExtension.java}{ScriptingExtension}; \source{core/src/main/java/info/openrocket/core/simulation/listeners/NewControlStepListener.java}{NewControlStepListener}; \source{core/src/main/java/info/openrocket/core/simulation/listeners/AirbrakesControllerListener.java}{AirbrakesControllerListener}.}

\section{Units, numeric primitives, and coordinate utilities}
The utility layer supplies the shared mathematics used by geometry, physics, plots and export. \class{Coordinate} stores Cartesian values and an optional weight; that weight can represent mass in one calculation and normal-force weighting in another. \class{WorldCoordinate} represents geographic location. \class{Quaternion} and \class{Transformation} encode orientation and rigid transforms, while inertia and bounding-box helpers support physical aggregation and geometry. These types do not remove the need to know the frame and units expected by the caller.

Interpolation, polynomial fitting, numerical statistics, array operations, noise generation and mathematical tolerance helpers are centralized in \code{util}. A linear interpolator over a thrust curve, a polynomial interpolation across aerodynamic regimes, and a pink-noise wind process solve different problems even though all are numerical helpers. Their boundary behavior, extrapolation rules and random-state ownership are part of the calling subsystem's contract.

The units package also does more than multiply by a display scale. Temperature conversions can include an offset. Caliber units depend on a rocket reference diameter, and percentage-of-length units depend on a relevant length. Fractional inch formatting and axis tick generation are presentation concerns. \class{FixedUnitGroup} gives custom quantities a fixed unit representation; it does not supply automatic physical dimension checking for arbitrary expressions. Values and comparators keep unit-aware display/sorting separate from internal SI calculations.

\sourceline{\source{core/src/main/java/info/openrocket/core/util/Coordinate.java}{Coordinate}; \source{core/src/main/java/info/openrocket/core/util/Quaternion.java}{Quaternion}; \source{core/src/main/java/info/openrocket/core/util/Transformation.java}{Transformation}; \source{core/src/main/java/info/openrocket/core/unit/UnitGroup.java}{UnitGroup}; \source{core/src/main/java/info/openrocket/core/unit/CaliberUnit.java}{caliber conversion}.}

\section{Change tracking, traversal, formatting, and object lifetimes}
Modification IDs identify changes that invalidate cached geometry or results. Listener lists and change-source helpers distribute notifications; invalidation helpers make disposed UI bindings stop acting on old models. Mutability guards distinguish objects still being assembled from completed data that should be read. Safety checks can detect inappropriate access, but they do not make a mutable document or engine safe for arbitrary concurrent writes.

Rocket visitors provide reusable traversals, including component and motor-mount listing and breadth/depth traversal strategies. They let consumers such as reports or selectors visit the component tree without embedding every traversal in a dialog. Active physical instance traversal remains a separate question from visiting each design node once.

The formatting package creates readable rocket/configuration descriptions and substitutes motor/configuration information into names. It is separate from expression evaluation and unit conversion. Configuration naming can depend on selected motors, while a persistent configuration ID remains the stable reference. Localization wraps translators, contextual keys and resource bundles so these labels can change with language without redefining stored identities.

File, JAR, reflection, platform, build-property and configuration helpers support startup and persistence. Their small size should not obscure their effect: resource lookup and class loading can determine whether a packaged application behaves like a development launch. The jar-in-jar bootstrap classes and platform helpers belong to this deployment boundary.

\section{Warnings, application logs, and network services}
Simulation warnings/errors are structured model output with types and associated context. They can survive in results and appear in reports. Application logging is a different channel: Swing logging classes integrate log buffers, filtering, levels and diagnostic displays with the logging backend. A debug message about a repaint is not a flight warning; a flight warning should not disappear merely because a console logging threshold changes.

Communication services retrieve welcome/release information and update metadata, and support bug-report workflows. Connection abstractions separate network access from response interpretation. MIT-specific update selection and replacement are described in Chapter~\ref{ch:build}; inherited update classes also remain in the tree. The existence of a network helper does not mean every simulation uses the network. Normal physics calculations operate on the already resolved design, motor and environment inputs.

ThrustCurve communication has its own request/response and curve decoding classes. A motor download must become a validated local motor record before it can participate in a configuration. Network metadata, local catalog entries and a runtime motor's ignition clock have separate lifetimes and failure modes.

\sourceline{\source{core/src/main/java/info/openrocket/core/logging/WarningSet.java}{structured warnings}; \source{core/src/main/java/info/openrocket/core/communication/ConnectionSource.java}{connection abstraction}; \source{core/src/main/java/info/openrocket/core/communication/WelcomeInfoRetriever.java}{welcome information}; \source{core/src/main/java/info/openrocket/core/formatting/RocketDescriptorImpl.java}{rocket description formatting}.}

\section{Inherited extension examples and developer utilities}
The extension example packages illustrate air starts, CSV saving, damping-moment changes, simulation printing, roll control and stopping conditions. Providers expose examples to discovery; runtime listeners implement their behavior. These examples are useful API references, but merely finding a class in the tree does not establish that a particular simulation has enabled it. The same distinction applies to Swing extension configurators: a panel edits configuration, while the core extension installs execution behavior.

The scripting support includes script-engine discovery and a GraalJS factory in addition to the simulation extension itself. Language availability and the packaged runtime therefore affect script execution. The separately bundled expression parser is for arithmetic/array expressions and does not replace the script-engine API or the MIT Java-board runtime.

Developer utility packages contain localization generators/reports, motor-correlation tools, optimizer exercises, preset conversion/editing/serialization utilities, and small desktop support tools. They are not all normal startup services. Build scripts, resources, tests and documentation complete the codebase around the runtime Java packages; Appendix~\ref{app:packages} provides an explicit coverage inventory so these less visible areas are not omitted from the architecture map.


\section{A maintenance checklist by boundary}
When changing a component, check geometry and instance placement, mass and inertia, aerodynamic behavior, renderer geometry, copy/identity semantics, serialization, and editor units. When changing an engine path, check event ordering, accepted-step semantics, stage branches, recovery-stepper transitions, cancellation, and result timestamps. When changing an FC task, check input acquisition, completion visibility, state-transition behavior, actuator ownership, recovery deduplication, and logging.

When adding a simulation option or channel, check old-file defaults, freshness signatures, ensemble behavior, CSV and ORK persistence, imported-data behavior, and plots. A new channel may need circular, held-state, or event semantics rather than ordinary scalar averaging. When adding concurrency, inspect mutable static state and external outputs in addition to object clones.

These checks follow the application's ownership boundaries. They are more useful than assuming that a successful GUI demonstration proves persistence, parallel execution, or the timing of a numerical callback.

\section{Practical investigation routes}
For a surprising trajectory, inspect configuration and launch options, then events, warnings, active stages, mass, aerodynamic coefficients and the selected stepper. For FC behavior, compare physical truth, measurements, task completion, command publication and actuator latching separately.

For differing runs, compare source/artifact identity, rocket/FC fingerprints, seeds, environment, configuration and extensions. For a save/reload mismatch, trace the relevant saver and loader fields before investigating a repaint.

For a visual mismatch, identify whether the disputed value belongs to design geometry, a renderer or recorded playback data. A changed drawing does not automatically change aerodynamic forces or an already recorded flight.
